\chapter{机械振动}
\setcounter{example_counter}{1}

\section{简谐振动的描述}

\subsection{简谐振动的表达式}

\subsubsection{位移}

\noindent
\begin{minipage}{0.73\textwidth}
    \setlength{\parskip}{0.5em}
    \setlength{\parindent}{2em}
    \linespread{1.3}
    \textbf{简谐振动}：离开平衡位置的位移$x$按余弦规律随时间变化。

    以平衡位置为基准，简谐振动位移的通式为
    
\end{minipage}
\hfill
\begin{minipage}{0.25\textwidth}
    \centering \includegraphics[width=\textwidth]{asset/ch5/T1.png}
\end{minipage}

\vspace{0.5em}

\begin{formula}
    $x=A\cos (\omega t+ \varphi)$
\end{formula}

\begin{itemize}
    \item $A$称为振幅，$\omega$称为角频率（或圆频率）
    \vspace{0.25em}
    \item 周期$T=\dfrac{2\pi}{\omega}$，频率$\nu=\dfrac{1}{T}=\dfrac{\omega}{2\pi}$
    \vspace{0.25em}
    \item $\omega t+ \varphi$称为相位，$t=0$时的相位$\varphi$称为初相位
\end{itemize}

\subsubsection{速度与加速度}

将$x$对$t$求导，可得速度$v$，再求导可得加速度$a$.

\begin{formula}
    $v=-A\omega \sin(\omega t+ \varphi),~a=-A\omega^2 \cos (\omega t+ \varphi)$
\end{formula}

从表达式可以看出，加速度与位移成正比且反向，这是简谐振动的运动学特征。

\begin{example}{简谐振动的运动学}
    某简谐振动的表达式为$x=A\cos ⁡(\omega t+\varphi)$ (SI)，已知$t=0$时的位移为$\SI{0.04}{m}$，速度为$\SI{0.12}{m/s}$，加速度为$\SI{-0.36}{m/s^2}$，求$A,\omega,\varphi$的值（$A,\omega>0,~|\varphi|<\pi$）。

    \tcblower

    将简谐振动表达式对$t$两次求导，得速度$v=-\omega A\sin{(\omega t+\varphi)}$，加速度$a=-A\omega^2 \cos (\omega t+ \varphi)$

    在$t=0$时，$x=A\cos\varphi=0.04,~v=-A\omega\sin\varphi=0.12,~a=-A\omega^2 \cos\varphi=0.36$
    
    两两相除得$\omega=3,~\tan\varphi=-1$，结合$\cos\varphi>0,~\sin\varphi<0$的要求，$\varphi=-\dfrac{\pi}{4}$符合题意。

    \vspace{0.5em}

    将$\cos\varphi=\dfrac{\sqrt{2}}{2}$代入$x=A\cos\varphi=0.04$，可得$A=0.04\sqrt{2}$
    
    \tcbline

    【补充练习】已知简谐振动的表达式为$x=0.1\cos{\left(8\pi t+\dfrac{2\pi}{3}\right)}$（SI），求最大速度、加速度，并求首次回到平衡位置的时间。
    
\end{example}


\subsection{旋转矢量法}

一个匀速圆周运动质点，其在$x$轴上的投影就是简谐振动，而且速度和加速度在$x$方向上的分量，也与简谐振动相同。利用这一关系，可以用逆时针旋转的矢量表示振动。

利用逆时针旋转、长度等于振幅$A$的矢量表示简谐振动，这种方法称为\textbf{旋转矢量法}。初始角度就是初相位$\varphi$，只画出初始矢量的位置，称为\textbf{相量图}，非常适合确定初相位。

\vspace{-0.5em}

\begin{figure}[htb]
    \centering
    \includegraphics[width=0.7\linewidth]{asset/ch5/T2.png}
\end{figure}

\vspace{-1em}

\begin{example}{简谐振动的表达式：旋转矢量法}
    已知简谐振动的振幅为$4$，周期为$2$. 当$t=0$时第一次通过$x=−2$处，且速度为正方向。求简谐振动的表达式（$A,\omega>0,~|\varphi|<\pi$），并求下一次通过$x=-2$位置的时间。

    \tcblower

    根据题意，$A=4,~T=2$，所以$\omega=\dfrac{2\pi}{T}=\pi$，简谐振动表达式为$x=4\cos{(\pi t+\varphi)}$

    如果直接将$t=0,~x=-2$代入表达式，会得到$\varphi=\pm\dfrac{2\pi}{3}$，还需根据速度方向进一步判断。

    \begin{minipage}{0.75\textwidth}
    
    而画出相量图，则可清楚看出初相位应该是第三象限角，应取$\varphi=-\dfrac{2\pi}{3}$

    可快速得出振动表达式为$x=4\cos{\left(\pi t-\dfrac{2\pi}{3}\right)}$

    根据向量图可知，下一次通过$x=-2$位置的相位是$\dfrac{2\pi}{3}$
    
    结合初相位，可知需要再等$\dfrac{4\pi}{3}$个相位。
    
    这相当于$\dfrac{4\pi/3}{2\pi}=\dfrac{2}{3}$个周期，因此再次通过$x=-2$的时间为$t=\dfrac{2}{3}T=\dfrac{4}{3}$

    \vspace{0.5em}
    
    \end{minipage}
    \hfill
    \begin{minipage}{0.22\textwidth}
        \vspace{-0.3em}
    
        \centering \includegraphics[width=\textwidth]{asset/ch5/T4.png}
    \end{minipage}

    \vspace{-1em}

    \tcbline

    【补充练习】已知简谐振动的图像如下图所示，求简谐振动的表达式。

    \vspace{-0.5em}

    \begin{center}
        \includegraphics[width=0.35\textwidth]{asset/ch5/T3.png}
    \end{center} 
    
\end{example}

\subsection{相位关系}

若两个振动的频率相同，则相位差$\Delta \varphi=\varphi_2−\varphi_1$不变，习惯上：
\begin{itemize}
    \item 若$\Delta\varphi=0$（或$2k\pi$），则称两个振动\textbf{同相}；若$\Delta\varphi=\pi$（或$(2k+1)\pi$），则称两个振动\textbf{反相}
    \item 若$\Delta\varphi>0$，一般说振动2\textbf{领先}$\Delta\varphi$的相位；若$\Delta\varphi<0$，一般说振动2\textbf{落后}$\Delta\varphi$的相位
\end{itemize}

\begin{figure}[htb]
    \centering
    \includegraphics[width=0.5\linewidth]{asset/ch5/T8.png}
\end{figure}

\begin{example}{简谐振动的表达式：相位关系}
    已知简谐振动的速度最大值为$4$，振幅为$2$，$t=0$时为速度最大值，且为正方向。
    
    (1)~求简谐振动的表达式。(2)~另一个振动与之反相，求其表达式。

    \tcblower
    
    \begin{minipage}{0.8\textwidth}
    
    (1)~根据题意，$A=2,~A\omega=4$，可得$\omega=2$

    根据$t=0$时取速度最大值，可画出右图所示的相量图，看出$\varphi=-\dfrac{\pi}{2}$

    因此简谐振动表达式为$x=2\cos{(2t-\dfrac{\pi}{2})}$
    
    \end{minipage}
    \hfill
    \begin{minipage}{0.18\textwidth}
    
        \centering \includegraphics[width=\textwidth]{asset/ch5/T6.png}
    \end{minipage}

    (2)~两个振动反相，则相位差为$\pi$，可得$\varphi_2=\dfrac{\pi}{2}$，因此表达式为$x=A\cos{(2t+\dfrac{\pi}{2})}$

    \tcbline

    【补充练习】两个质点做同振幅、同频率的简谐振动，在振动过程中，它们总在经过振幅一半的地方相遇，且运动方向相反。求他们的相位差。
    
\end{example}

\section{简谐振动的力与能量}

\subsection{简谐振动的动力学}

根据简谐运动的表达式，可知$a=-\omega^2 x$，加速度与位移正比且反向，这是简谐振动的的动力学特征。

结合牛顿第二定律，可知简谐振动的回复力与位移正比且反向。
\begin{formula}
    $F=-kx~~~\text{或}~~~m\dfrac{dx^2}{dt^2}=-kx$
\end{formula}

结合牛顿第二定律，$F=ma$，其中$a=-\omega^2 x$，可得简谐振动的角频率和周期为
\begin{formula}
    $\omega=\sqrt{\dfrac{k}{m}},~~T=2\pi\sqrt{\dfrac{m}{k}}$
\end{formula}

\subsection{简谐振动的能量}

\vspace{-1em}

\noindent
\begin{minipage}{0.78\textwidth}
    \setlength{\parskip}{0.5em}
    \setlength{\parindent}{2em}
    \linespread{1.3}
    在简谐振动的过程中，动能和势能相互转化，转化的过程中总能量不变。

    \begin{itemize}
        \item 在平衡位置，动能最大而势能为零
        \item 在位移最大处，动能为零而势能最大
    \end{itemize}
    
\end{minipage}
\hfill
\begin{minipage}{0.2\textwidth}
    \centering \includegraphics[width=\textwidth]{asset/ch5/T18.png}
\end{minipage}

%$E_k=\dfrac{1}{2}mv^2=\dfrac{1}{2}mA^2\omega^2\sin^2(\omega t+\varphi),~E_p=\dfrac{1}{2}kx^2=\dfrac{1}{2}kA^2\cos^2(\omega t+\varphi)$，结合$\omega=\sqrt{\dfrac{k}{m}}$，可得

%$E=E_k+E_p=\dfrac{1}{2}kA^2=\dfrac{1}{2}m\omega^2 A^2$

\begin{formula}
    $E=E_k+E_p=\dfrac{1}{2}mv^2+\dfrac{1}{2}kx^2=\dfrac{1}{2}kA^2=\dfrac{1}{2}m\omega^2 A^2$
\end{formula}

\begin{example}{简谐振动的能量}
    有一弹簧振子，弹簧劲度系数为$k=\SI{25}{N/m}$，振子振动的初动能为$\SI{2}{J}$，初势能为$\SI{6}{J}$.
    
    求：(1)~振幅；(2)~当位移是多大时，势能与动能相等。

    \tcblower

    

    (1)~总能量$E=E_k+E_p=\SI{8}{J}$，当振幅最大时，全为势能，$E=\dfrac{1}{2}kA^2$，得$A=\SI{0.8}{m}$

    (2)~势能与动能相等时，$E_p=\dfrac{1}{2}E=\SI{4}{J}$，利用$E_p=\dfrac{1}{2}kx^2$，得$x=0.4\sqrt{2}\unit{m}$

    \tcbline

    【补充练习】某弹簧振子重物的质量为$\SI{0.01}{kg}$，振动表达式$x=0.1\cos\left(8\pi t+\dfrac{2\pi}{3}\right)$（SI）.~求：(1)~弹簧的劲度系数；(2)~通过平衡位置时的速率；(3)~动能等于势能3倍时的位移。
    
\end{example}

\subsection{简谐振动的实例}

\subsubsection{弹簧振子}

弹簧振子的$A,\varphi$由初始条件确定，但角频率和周期只由系统本身决定。

\noindent
\begin{minipage}{0.73\textwidth}
    \setlength{\parskip}{0.5em}
    \setlength{\parindent}{2em}
    \linespread{1.3}
    弹簧振子的回复力是弹簧的弹力，满足$F=-kx$，所以是简谐振动。

    
    
\end{minipage}
\hfill
\begin{minipage}{0.25\textwidth}
    \centering \includegraphics[width=\textwidth]{asset/ch5/T1.png}
\end{minipage}

\begin{formula}
    $\dfrac{d^2x}{dt^2}=-\dfrac{k}{m}x,~~\omega=\sqrt{\dfrac{k}{m}},~~T=\dfrac{2\pi}{\omega}=2\pi\sqrt{\dfrac{m}{k}}$
\end{formula}

\note{若弹簧振子竖直放置，以平衡位置为基准，重力势能与弹性势能之和仍为$\dfrac{1}{2}kx^2$，所有结论仍成立。}

\begin{example}{弹簧振子}
    \begin{minipage}{0.73\textwidth}
        光滑水平面上的弹簧振子有质量为$M$的木块和劲度系数为$k$的弹簧构成。弹簧处于原长，一个质量为$m$的子弹以速度$u$射入木块并留在其中，系统开始振动。求振动方程（向右为正）。 
    \end{minipage}
    \hfill
    \begin{minipage}{0.25\textwidth}
        \centering \includegraphics[width=\textwidth]{asset/ch5/T10.png}
    \end{minipage}
    
    \tcblower

    

    在射入后瞬间，由动量守恒，$mu=(m+M)v$，此时系统能量为$E=E_k=\dfrac{1}{2}(M+m)v^2$

    振动过程中总能量不变，因此在振幅最大处，$E=E_p=\dfrac{1}{2}kA^2$，解得$A=\dfrac{mu}{\sqrt{k(M+m)}}$

    利用简谐运动的动力学公式，可得$\omega=\sqrt{\dfrac{k}{M+m}}$，利用旋转矢量法可知初相位$\varphi=-\dfrac{\pi}{2}$
    
    因此振动方程为$x=\dfrac{mu}{\sqrt{k(M+m)}}\cos\left(\sqrt{\dfrac{k}{M+m}}t-\dfrac{\pi}{2}\right)$

    \tcbline

    \vspace{-1em}
    
    \begin{minipage}{0.88\textwidth}
        【补充练习】劲度系数$\SI{10}{N/m}$的轻弹簧下端挂有质量为$\SI{1}{kg}$的盘子。现有质量也为$\SI{1}{kg}$的物体在高为$\SI{1}{m}$处自由下落，并与盘子粘在一起，系统开始振动。求振动方程。
    \end{minipage}
    \hfill
    \begin{minipage}{0.1\textwidth}
        \centering \includegraphics[width=\textwidth]{asset/ch5/T11.png}
    \end{minipage}
    
\end{example}

\subsubsection{单摆}

\noindent
\begin{minipage}{0.85\textwidth}
    \setlength{\parskip}{0.5em}
    \setlength{\parindent}{2em}
    \linespread{1.3}
    单摆的回复力是重力在切线方向的分力$mg\sin \theta$，若摆动幅度很小，可用小角近似$\sin \theta \sim \theta$，因此可认为回复力$F=-mg\theta$，与$\theta$正比且反向。这说明在摆角很小的情况下，单摆的运动是简谐振动。

    摆球的切向加速度$a_t=l\alpha=l\dfrac{d^2\theta}{dt^2}$，所以动力学方程为$ml\dfrac{d^2\theta}{dt^2}=-mg\theta$，据此可得

    \begin{formula}
        $\dfrac{d^2\theta}{dt^2}=-\dfrac{g}{l}\theta,~~\omega=\sqrt{\dfrac{g}{l}},~~T=2\pi\sqrt{\dfrac{l}{g}}$
    \end{formula}
    
\end{minipage}
\hfill
\begin{minipage}{0.12\textwidth}
    \centering \includegraphics[width=\textwidth]{asset/ch5/T12.png}
\end{minipage}

注意不要把单摆振动的相位与摆角混淆。
\begin{itemize}
    \item 最大摆角对应最大位移，相位为0和$\pi$
    \item 摆角为0对应平衡位置，相位为$、pm\dfrac{\pi}{2}$
\end{itemize}

\subsubsection{复摆}

\noindent
\begin{minipage}{0.83\textwidth}
    \setlength{\parskip}{0.5em}
    \setlength{\parindent}{2em}
    \linespread{1.3}
    绕不过质心的水平固定轴转动的刚体，称为复摆，需要用转动定律进行分析。

    在摆角很小的情况下，回复力矩$M=-mgl\sin\theta\sim -mgl\theta$，其中$l$表示质心到转轴的距离。回复力矩与$\theta$正比且反向，说明复摆的运动是简谐振动。
    
    利用转动定律，可列出振动的动力学方程$J\dfrac{d^2\theta}{dt^2}=-mgl\theta$，据此可得
    
\end{minipage}
\hfill
\begin{minipage}{0.15\textwidth}
    \centering \includegraphics[width=\textwidth]{asset/ch5/T13.png}
\end{minipage}

\begin{formula}
    $\dfrac{d^2\theta}{dt^2}=-\dfrac{mgl}{J}\theta,~~\omega=\sqrt{\dfrac{mgl}{J}},~T=\dfrac{2\pi}{\omega}=2\pi\sqrt{\dfrac{J}{mgl}}$
\end{formula}

\begin{example}{单摆与复摆}

    \begin{minipage}{0.85\textwidth}
        一均匀细圆环质量为$m$，半径为$R$，挂在墙上的钉子上。求微小摆动的周期。

        提示：圆环相对于钉子的转动惯量$J=2mR^2$.

    \end{minipage}
    \hfill
    \begin{minipage}{0.1\textwidth}
        \centering \includegraphics[width=\textwidth]{asset/ch5/T19.png}
    \end{minipage}

    \tcblower
    
    圆环若偏离很小的$\theta$角度，则回复力矩$M=-mgR\sin\theta \approx -mgR\theta$

    \vspace{0.5em}

    振动方程为$J\dfrac{d^2\theta}{dt^2}=-mgR\theta$，整理得$\dfrac{d^2\theta}{dt^2}=-\dfrac{g}{2R}\theta$.~因此$\omega=\sqrt{\dfrac{g}{2R}}$，$T=\dfrac{2\pi}{\omega}=2\pi \sqrt{\dfrac{2R}{g}}$
    \tcbline
    【补充练习】将一个单摆的摆球向负方向拉开很小的$\theta$角并放手，使之自由运动。它的初相位是多少？摆球首次来到平衡位置时的相位是多少？
\end{example}

\section{同方向简谐振动的合成}

\subsection{同方向同频率}

\noindent
\begin{minipage}{0.73\textwidth}
    \setlength{\parskip}{0.5em}
    \setlength{\parindent}{2em}
    \linespread{1.3}
    同方向、同频率的两个振动$x_1=A_1\cos{(\omega t+\varphi_1)},~x_2=A_2\cos{(\omega t+\varphi_2)}$，合成的结果仍然是同频率的简谐振动，参数可根据旋转矢量法进行确定。

    \begin{itemize}
    \item 若两个振动同相（$Δ\varphi=0$），则振幅最大，$A=A_1+A_2$
    \item 若两个振动反相（$Δ\varphi=\pi$），则振幅最小，$A=|A_1-A_2|$
    \item 其余情况，振幅$A$介于二者之间
\end{itemize}
\end{minipage}
\hfill
\begin{minipage}{0.25\textwidth}
    \centering \includegraphics[width=\textwidth]{asset/ch5/T14.png}
\end{minipage}

\begin{example}{同方向同频率的振动合成}
    有两个同方向的振动$x_1=4\cos⁡{(2t+\dfrac{\pi}{4})},~x_2=4\cos{⁡(2t+\dfrac{3\pi}{4})}$进行合成。 (1) 求合振动的方程；(2) 若还有第三个振动$x_3=3\cos⁡(2t+\varphi)$也参与合成，$\varphi$取何值时，合振动的振幅取得最大值和最小值，最大、最小值分别是多少？

    \tcbline

    %{\fangsong 同方向同频率振动合成的一般方法是旋转矢量法。}

    %\tcbline

    \begin{minipage}{0.76\textwidth}
        (1)~画出相量图可看出，合振动的振幅为$4\sqrt{2}$，初相位为$\dfrac{\pi}{2}$，频率不变。
    
        因此合振动方程为$x=4\sqrt{2}\cos{(2t+\dfrac{\pi}{2}})$
    
        (2)~要使合振动的振幅最大，应使振动同相，$\varphi=\dfrac{\pi}{2}$，最大振幅$A=4+3=7$

        \vspace{0.3em}
    
        要使合振动的振幅最小，应使振动反相，$\varphi=-\dfrac{\pi}{2}$，最大振幅$A=4-3=1$
    \end{minipage}
    \hfill
    \begin{minipage}{0.22\textwidth}
        \centering \includegraphics[width=\textwidth]{asset/ch5/T15.png}
    \end{minipage}

    \tcblower
    【补充练习】设有两个同方向简谐振动$x_1=A\cos(\omega t+\dfrac{\pi}{4}),~x_2=\sqrt{3}A\cos(\omega t+\dfrac{3\pi}{4})$. (1)~求合振动的表达式；(2)~再与第三个振动合成，使合振动的振幅达到$4A$，并让第三个振动的振幅尽可能小。

\end{example}

\subsection{同方向不同频率}
若两个振动的频率略有差异，合振动的振幅出现周期性变化，称为\textbf{拍}，拍频就是两个振动的频率之差。

\begin{formula}
    $\nu=|\nu_2-\nu_1|=\left|\dfrac{\omega_1}{2\pi}-\dfrac{\omega_2}{2\pi}\right|$
\end{formula}

%\begin{figure}[htb]
%    \centering
%    \includegraphics[width=0.35\linewidth]{asset/ch5/T16.png}
%\end{figure}

\begin{example}{拍现象}
    音叉A与频率$\SI{415}{Hz}$的音叉B同时振动时，可在每秒内听到$2$拍。则音叉A与频率$\SI{420}{Hz}$的音叉C同时振动时，每秒内能听到几拍？
    \tcbline
    根据音叉A、B的振动可知，$|\nu_A-\nu_B|=\SI{2}{Hz}$，得$\nu_A=\SI{417}{Hz}$或$\SI{413}{Hz}$.

    因此音叉A、C振动时，拍频$\nu^\prime=|\nu_A-\nu_B|=\SI{3}{Hz}$或$\SI{7}{Hz}$.
    \tcbline
    【补充练习】一个$\SI{256}{Hz}$的音叉和待测音叉同时鸣响，若在$\SI{20}{s}$内听到$10$拍，求待测音叉的频率。
\end{example}

\section{拓展内容}

\subsection{简谐振动的判定}

要想判定一个运动是简谐振动，有以下三种等价的方法：
\begin{itemize}
    \item 受力特征：回复力与位移正比且反向，$F=-kx$
    \item 振动方程：$\dfrac{d^2x}{dt^2}+\omega^2x=0$（即$a=-\omega^2 x$）
    \item 能量特征：总能量不变，且势能$E_p=\dfrac{1}{2}kx^2$
\end{itemize}

\begin{example}{简谐振动的判定*}
    
    
    \begin{minipage}{0.8\textwidth}
    
    长为$l$、横截面积为$S$的木块浮在水面上，静止时浸没深度为$h$。将其稍向下压后松手，证明其之后的运动是简谐振动，并求振动周期。
    
    
    
    \end{minipage}
    \hfill
    \begin{minipage}{0.18\textwidth}
    
        \centering \includegraphics[width=\textwidth]{asset/ch5/T7.png}
    \end{minipage}

    

    \tcblower

    在平衡位置，浮力等于重力，$\rho_{\text{水}} gSh=\rho_{\text{木}}Slg$，得$\rho_{\text{水}}h=\rho_{\text{木}}l$

    以平衡位置为原点，向下为正。若向下压$x$的距离，则合力$F$等于重力减浮力。

    $F=\rho_{\text{木}}Slg-\rho_{\text{水}} gS(h+x)=-\rho_{\text{水}}gSx$，这说明回复力与位移成正比且反向，木块做简谐振动。

    对照简谐振动的动力学表达式，$k=\rho_{\text{水}}gS$，木块$m=\rho_{\text{木}}Sl$

    \vspace{0.5em}

    代入公式可得$\omega=\sqrt{\dfrac{k}{m}}=\sqrt{\dfrac{\rho_{\text{水}}gS}{\rho_{\text{木}}Sl}}=\sqrt{\dfrac{g}{h}}$，周期$T=\dfrac{2\pi}{\omega}=2\pi \sqrt{\dfrac{h}{g}}$
        
    \tcbline 

    \begin{minipage}{0.82\textwidth}
    
    【补充练习】横截面均匀的光滑U形管中有适量液体，液体的总长度为$L$，忽略摩擦和液体黏性。证明液面的小幅上下振动是简谐振动，并求出振动周期。
    
    \end{minipage}
    \hfill
    \begin{minipage}{0.15\textwidth}
        \vspace{-2em}
        \centering 
\includegraphics[width=\textwidth]{asset/ch5/T9.png}
    \end{minipage}
    

\end{example}

\subsection{阻尼振动与受迫振动}

\subsubsection{阻尼振动}

设自振频率$\omega_0$的系统在振动过程中，受到$f=-\gamma v$的阻力，记阻尼系数$\beta=\dfrac{\gamma}{2m}$

\begin{itemize}
    \item 若$\beta>\omega_0$，称为\textbf{过阻尼}，不能振动
    \item 若$\beta=\omega_0$，称为\textbf{临界阻尼}，恰好不能振动
    \item 若$\beta<\omega_0$，称为\textbf{欠阻尼}，可以振动
\end{itemize}

在欠阻尼的情况下，振动周期变长$T=\dfrac{2\pi}{\sqrt{\omega_0^2-\beta^2}}$，振幅和能量逐渐减小，$A=A_0 e^{-\beta t},~E=E_0 e^{-2\beta t}$.

振幅减小到$\dfrac{1}{e}$的时间称为时间常量（鸣响时间）$\tau=\dfrac{1}{2\beta}$

在鸣响时间内振动的次数的$2\pi$倍称为品质因数（Q值），$Q=2\pi\dfrac{\tau}{T}=\omega \tau$

\begin{example}{阻尼振动*}
    一单摆在空气中摆动，初始振幅为$5^\circ$.~若经过$\SI{100}{s}$，振幅减为$4^\circ$.~求：再经过多长时间，振幅减为$2^\circ$，并求出鸣响时间。
    \tcblower
    根据$A=A_0 e^{-\beta t}$，可得$t=\dfrac{1}{\beta}\ln{\dfrac{A_0}{A}}$.~代入$A_0=5^\circ,~A=4^\circ,~t=\SI{100}{s}$，可求出$\beta\approx \SI{2.2e-3}{s^{-1}}$
    
    再代入$A_0=4^\circ,~A=2^\circ$的数据，可得$t\approx \SI{310.7}{s}$.~进而鸣响时间$\tau=\dfrac{1}{2\beta}\approx \SI{448}{s}$
    \tcbline
    【补充练习】阻尼振动的起始振幅为$\SI{3.0}{cm}$，$\SI{10}{s}$后变为$\SI{1.0}{cm}$，再经过多长时间振幅变为$\SI{0.3}{cm}$.
\end{example}

\subsubsection{受迫振动}

在周期性驱动力的作用下，系统振动的频率等于驱动力的频率，而不是自振频率。

驱动力的频率与自振频率越接近，振幅越大。当二者相等时，振幅都达到最大值，称为\textbf{共振}现象。

\subsection{相互垂直的简谐振动合成}

%\subsubsection{相互垂直的简谐振动合成}

\vspace{-0.5em}

\noindent
\begin{minipage}{0.65\textwidth}
    \setlength{\parskip}{0.5em}
    \setlength{\parindent}{2em}
    \linespread{1.3}
    若质点在$x,y$方向的运动都是简谐运动，且频率相等，则合运动的轨迹是直线或椭圆，取决于相位差。
    
    \begin{itemize}
    \item 若相位差为$0$或$\pi$，则轨迹为直线
    \item 若$y$的相位比$x$领先，则为右旋的椭圆（顺时针转动）
    \item 若$y$的相位比$x$落后，则为左旋的椭圆（逆时针转动）
\end{itemize}
    
\end{minipage}
\hfill
\begin{minipage}{0.33\textwidth}
    \centering \includegraphics[width=\textwidth]{asset/ch5/T20.png}
\end{minipage}

\vspace{-0.5em}

\noindent
\begin{minipage}{0.78\textwidth}
    \setlength{\parskip}{0.5em}
    \setlength{\parindent}{2em}
    \linespread{1.3}
    若质点在$x,y$方向的运动都是简谐运动，且频率是简单的整数比，则质点合运动的轨迹是一个闭合曲线，称为\textbf{李萨如图形}。

    沿着闭合曲线，两个方向振动的频率之比，等于到达$x,y$最大位移的次数之比。若到达角点，则各算半次。
    
\end{minipage}
\hfill
\begin{minipage}{0.2\textwidth}
    \centering \includegraphics[width=\textwidth]{asset/ch5/T17.png}
\end{minipage}

%\begin{formula}
%    $\dfrac{\nu_x}{\nu_y}=\dfrac{x\text{方向最大位移次数}}{y\text{方向最大位移次数}}$
%\end{formula}

\note{若$x,y$方向的运动不成比例，则质点的运动轨迹不闭合，但很容易锁定到简单整数比。}

\begin{example}{李萨如图形*}
    \vspace{-0.5em}
    \begin{minipage}{0.88\textwidth}
        示波器的水平和垂直输入端分别加上余弦式交变电压，荧光屏上出现如图所示的曲线。已知水平方向振动频率为$\SI{2.7e4}{Hz}$，求垂直方向的振动频率。
        
    \end{minipage}
    \hfill
    \begin{minipage}{0.08\textwidth}
        \centering \includegraphics[width=\textwidth]{asset/ch5/T21.png}
    \end{minipage}
    \vspace{-0.5em}
    \tcblower

    \linespread{1.3}
    
    假想沿着曲线行走一圈，可以到达6次$x$最大位移和4次$y$最大位移，这说明在相同时间内，$x$方向振动了3个周期，而$y$方向只振动了2个周期，所以$\dfrac{\nu_x}{\nu_y}=\dfrac{6}{4}$，可得$\nu_y=\SI{1.8e4}{Hz}$
    
    \tcbline
    【补充练习】写出下列李萨如图形两个方向振动的频率之比。

    \begin{center}
        \includegraphics[width=0.6\textwidth]{asset/ch5/T22.png}
    \end{center}
    
\end{example}